%% Cover letter for the CVIU submission (portal: required file "Cover letter", uploaded as PDF).
%% CVIU's guide for authors sets no content requirements for the cover letter.
%% Layout and fixed fields only: the body is Dean's to write against the content points (red \point markers).
%% Agent builds: ./build.sh -> build_agent/ (never in this directory; the IDE compiles here on save).
\documentclass[11pt,a4paper]{article}
\usepackage[T1]{fontenc}
\usepackage{newtxtext,newtxmath}
\usepackage[a4paper,margin=25mm]{geometry}
\usepackage{xcolor}
\usepackage{url}
\usepackage[british]{babel}
\setlength{\parindent}{0pt}
\setlength{\parskip}{0.8em}
\pagestyle{empty}
%% Content-point marker: prints in red so an unfinished letter is obvious. Remove every \point before upload.
\newcommand{\point}[1]{\textcolor{red}{[#1]}}

\begin{document}

%% Letterhead (postal address: use the same one as the manuscript's affiliation once confirmed)
{\large\bfseries Dean Mills}\\
Department of Electrical, Electronic and Computer Engineering\\
University of Pretoria, Pretoria, South Africa\\
\url{deansmills@gmail.com}/\url{u13081650@tuks.co.za}

\vspace{1em}
\today

\vspace{1em}
%% Name of the Editor-in-Chief from the journal's editorial-board page, if wanted
The Editor-in-Chief\\
\textit{Computer Vision and Image Understanding}

\vspace{1em}
\textbf{Re: Submission of the manuscript ``Detector-independent keypoint grouping for multi-person pose estimation with a graph attention network''}

Dear Editor,

On behalf of my supervisor, I submit the manuscript ``Detector-independent keypoint grouping for multi-person pose estimation with a graph attention network'' by Dean Mills and Hans Grobler for consideration as an original research article in \textit{Computer Vision and Image Understanding}.

Bottom-up multi-person pose estimators assign detected keypoints to people using a grouping mechanism trained jointly with the detector, so the grouping cannot be reused once the detector changes.
The manuscript treats grouping as a separate module and measures the cost of that separation.
The proposed grouping module, PG-GAT, is a graph attention network that embeds keypoints so that $k$-means recovers the people.
Trained on synthetic scenes and fine-tuned on COCO, it is applied to HigherHRNet detections and, without retraining, to YOLO11x-pose and RTMO-L; the autonomous pipeline trails the native grouping of all three detectors by 0.020 to 0.022 in pose grouping accuracy; at the system level, the cost is 0.048 average precision on identical keypoints, rising to 0.076 in scenes with six or more people.
Learned grouping heads and modularity-based clustering did not improve on $k$-means, whereas embedding changes produced gains that transfer to real data, identifying the embedding as the limiting component.

The work falls within the journal's research area of matching and recognition and follows recent keypoint and skeleton work published in the journal (Yin et al., 2024; Tanaka et al., 2024; Simoni et al., 2025; Clinckemaillie et al., 2026).

The manuscript is original, has not been published previously, and is not under consideration for publication elsewhere.
Both authors have read and approved the submitted version.
The manuscript builds on my MEng dissertation at the University of Pretoria, which is currently under examination; an academic thesis does not constitute prior publication under Elsevier's policy.

The authors declare no competing interests, and the use of generative AI in preparing the manuscript is declared in the manuscript.
The code, the trained checkpoints and the synthetic dataset can be made publicly available once the examination of the dissertation is complete; until then, they can be made available to the editor and reviewers on request.

Thank you for considering the manuscript.

Yours sincerely,

\vspace{2.5em}
Dean Mills\\
Corresponding author, on behalf of both authors

\end{document}
